\chapter{Conclusion}
\label{ch:conclusion}

The four theorems of this monograph were introduced, in
Chapter~\ref{ch:introduction}, as four questions one can ask of a single
admissibility relation. Having proven each in turn, it is worth stating
plainly what has and has not been established.

Representation Invariance (Theorem~\ref{thm:invariance}) established
that sameness of representation, properly understood, is sameness of
admissibility structure and nothing else; substrate is invisible to the
theorem except insofar as it constrains $R$. Admissibility Preservation
(Theorem~\ref{thm:preservation}) established that coherence of a
trajectory is not a separate property requiring its own definition, but
exactly the nondegeneracy of the admissibility relation examined along
that trajectory. Monotonic Relaxation (Theorem~\ref{thm:relaxation})
established that a scalar measure of inadmissibility decreases along
repair trajectories specifically, as a direct consequence of how repair
is already characterized elsewhere in this program, and Corollary
\ref{cor:relaxation-continuation} fixed the precise, non-automatic
relationship between this claim and the unrelated ordering relation of
monotonic continuation. Non-Identifiability
(Theorem~\ref{thm:nonident-structural}) established that failures of
historical reconstruction are not contingent misfortunes but a
structural guarantee whenever observation is constrained to reachable-set
structure and the space of possible histories exceeds what that
structure can encode, and Corollary~\ref{cor:nonident-diagnosis} gave a
principled way to diagnose, in any particular case, whether a
reconstruction failure reflects a correctable error or a genuine limit.

Two things follow from having stated these four theorems together rather
than separately, as they have variously appeared in earlier work in this
program. The first is economy: each theorem's proof drew only on
definitions already established elsewhere, and none required
introducing new primitive vocabulary. The second, and more important, is
that the four theorems are visibly not four independent facts about four
different subjects --- representation, coherence, repair, and
reconstruction --- but four consequences of one relation, $x \adm{R} y$,
examined respectively at a single state, along a trajectory, under a
restorative dynamics, and under a lossy projection. What varies across
the four chapters is not the object of study but the operation performed
on it. This is, in the end, the thesis this monograph exists to make
precise: that admissibility is not one tool among several available for
studying structure under transformation, but the common substrate from
which invariance, coherence, repair, and the limits of recovery are all
derivable in turn.
